% =============================================================================
% Dokumentklasse und Spracheinstellung
% =============================================================================
% für elektronische Veröffentlichung (Standard)
\documentclass[11pt,oneside,paper=A4,DIV=11,BCOR=0mm,abstract=true,headsepline,headings=normal,ngerman]{scrreprt}

\usepackage[ngerman]{babel}
\usepackage[utf8]{inputenc}
\usepackage[T1]{fontenc}

% =============================================================================
% Schriftart & Typographie
% =============================================================================
\usepackage{libertine}
\usepackage{libertinust1math}
\usepackage[auto]{microtype}
\usepackage[htt]{hyphenat}          % Silbentrennung in \texttt
\setlength{\emergencystretch}{3em}  % Notfall-Streckung bei langen Code-Snippets

% Verhindert "Schusterjungen" und "Hurenkinder" (einzelne Zeilen auf neuer Seite)
\clubpenalty = 10000
\widowpenalty = 10000
\displaywidowpenalty = 10000

% Floats auf Float-Seiten oben ausrichten statt vertikal zentrieren
\makeatletter
\setlength{\@fptop}{0pt}
\setlength{\@fpbot}{0pt plus 1fil}
\makeatother

% =============================================================================
% Mathe, Symbole, Einheiten
% =============================================================================
\usepackage{amsmath}
\usepackage{amsxtra}
\usepackage{eurosym}
\usepackage{siunitx}  
\sisetup{locale=DE}
\usepackage[version=4]{mhchem}
\usepackage{chemfig}

% =============================================================================
% Bilder, Tabellen, Quellcode
% =============================================================================
\usepackage{graphicx}
\setkeys{Gin}{width=\textwidth,height=0.45\textheight,keepaspectratio}
% Abbildung mit benutzerdefinierter Breite (Höhenlimit bleibt einheitlich)
\newcommand{\narrowfig}[2][0.8]{%
  \includegraphics[width=#1\textwidth,height=0.45\textheight,keepaspectratio]{#2}}
% Für gestapelte Subfigures: kleineres Höhenlimit damit beide passen
\newcommand{\stackedfig}[1]{%
  \includegraphics[width=\textwidth,height=0.35\textheight,keepaspectratio]{#1}}
\usepackage{float}
\usepackage{multirow,multicol,booktabs}
\usepackage{threeparttable}
\usepackage{longtable}
\usepackage{rotating}
\usepackage{pdflscape}  % Landscape-Seiten für Anhang-Tabelle
\usepackage{ltablex} % Kombination aus tabularx und longtable
\usepackage{subcaption}
\captionsetup[subtable]{position=top}
\usepackage{pdfpages}
\usepackage{listings}
\lstdefinelanguage{TypeScript}{
  keywords={break, case, catch, continue, debugger, default, delete, do, else,
    finally, for, function, if, in, instanceof, new, return, switch, this,
    throw, try, typeof, var, void, while, with, const, let, of, class,
    extends, import, export, from, interface, type, enum, implements,
    abstract, as, async, await, declare, keyof, namespace, readonly,
    undefined, null, true, false, any, string, number, boolean, never, void,
    object, symbol, bigint, unknown},
  sensitive=true,
  comment=[l]{//},
  morecomment=[s]{/*}{*/},
  morestring=[b]',
  morestring=[b]",
  morestring=[b]`,
}
\usepackage{blindtext}
\usepackage{csquotes}

% Darstellung von URL (wird von hyperref später erweitert)
\usepackage{url}
\urlstyle{same}

% \code{...}: Inline-Code mit Zeilenumbruch an Sonderzeichen
% Für reine Texte ohne LaTeX-Befehle; \texttt bleibt für Argumente mit \, etc.
\newcommand{\code}[1]{\texttt{\nolinkurl{#1}}}

% Fussnoten, auch für Tabellen
\usepackage{footnote}
\makesavenoteenv{tabular} 

% =============================================================================
% Literaturverzeichnis (BibLaTeX)
% =============================================================================
\usepackage[
backend=biber,
style=numeric,
sorting=none,
citestyle=numeric-comp,
maxbibnames=10,
giveninits=true,
]{biblatex}

\renewcommand*{\labelnamepunct}{\addcolon\addspace}

% Pfad zur .bib Datei
\addbibresource{literatur/literaturdatenbank.bib}

% =============================================================================
% Nomenklatur / Abkürzungsverzeichnis
% =============================================================================
\usepackage[intoc, german]{nomencl}
\makenomenclature

\RedeclareSectionCommand[tocbeforeskip=0.5em plus 0.3em]{chapter}
\renewcommand{\nomname}{Abkürzungs- und Symbolverzeichnis}
\setlength{\nomitemsep}{-\parsep}
\newlength{\nomgroupstartsep}
\setlength{\nomgroupstartsep}{18pt}

\newcommand{\nomenclheader}[1]{\item[\hspace*{-\itemindent}\itshape#1]{\ } \vspace{6pt}\par}

\renewcommand\nomgroup[1]{%
	\itemsep\nomgroupstartsep
	\ifstrequal{#1}{A}{\nomenclheader{Abkürzungen}}{%
		\ifstrequal{#1}{E}{\newpage\nomenclheader{Lateinische Symbole}}{%
			\ifstrequal{#1}{G}{\nomenclheader{Griechische Symbole}}{%
				\ifstrequal{#1}{O}{\nomenclheader{Operatoren}}{%
					\ifstrequal{#1}{X}{\nomenclheader{Tiefgestelle Indizes}}{%
						\ifstrequal{#1}{Y}{\nomenclheader{Hochgestelle Indizes}}{}}}}}}%
	\itemsep\nomitemsep
}

% =============================================================================
% Fußnoten-Design
% =============================================================================
% Linksbündige Fußnotenmarkierungen
\deffootnote{1.5em}{1em}{%
	\makebox[1.5em][l]{\thefootnotemark}%
}

% Fußnoten nicht über Seiten umbrechen
\interfootnotelinepenalty=10000

% =============================================================================
% Beschriftungen (Captions) und KOMA-Fonts
% =============================================================================
\addtokomafont{caption}{\small}
\setkomafont{captionlabel}{\sffamily\bfseries}
\setkomafont{author}{\large}
\setkomafont{date}{\large}
\setkomafont{publishers}{\large}
\setkomafont{dedication}{\normalsize}
\addtokomafont{dedication}{\raggedright}

\renewcaptionname{ngerman}{\figurename}{Abb.}
\renewcaptionname{ngerman}{\tablename}{Tab.} 

% Eigene Tabellenumgebungen
\newenvironment{tabular10}{%
	\fontsize{10}{12}\selectfont\tabular
}{%
	\endtabular
}

\newenvironment{tabular7}{%
	\fontsize{7}{12}\selectfont\tabular
}{%
	\endtabular
}

% =============================================================================
% Datenimport & Hyperref (Muss fast am Ende stehen!)
% =============================================================================
%\def\VARtyp{Bachelorarbeit}
\def\VARtyp{Bachelorarbeit}

\def\VARtitel{Performance- und Ressourcenanalyse von
	React und Svelte bei steigender UI-Dichte und
	Update-Frequenz}


\def\VARname{Jonas Öhler}

\def\VARmatrikelnummer{44338}

\def\VARstudiengang{Medieninformatik}

\def\VARhochschule{Hochschule der Medien Stuttgart}

\def\VARabgabedatum{20.07.2026}

\def\VARabschlussbezeichnung{Bachelor of Science}

\def\VARerstbetreuung{Prof. Dr. Simon Wiest}
\def\VARzweitbetreuung{Prof. Dr. Tobias Jordine}

%wenn keine externe Arbeit, dann extern auf false setzen



\usepackage[
pdftitle={\VARtitel},
pdfsubject={\VARtyp},
pdfauthor={\VARname},
pdfkeywords={},  
hidelinks,      % Links im PDF nicht umranden
breaklinks=true % Zeilenumbruch bei langen URLs erlauben
]{hyperref}

% Cleveref als Abschluss für intelligente Querverweise
\usepackage[german]{cleveref}

% Umbruchzeichen für \code{...} nach hyperref erweitern
\expandafter\def\expandafter\UrlBreaks\expandafter{\UrlBreaks%
  \do\(\do\)\do\,\do\;\do\_\do\=\do\.\do\/\do\@\do\'}